\documentclass[11pt]{article}
\usepackage[margin=1in]{geometry}
\usepackage{enumitem}
% =============================================================================
% Preambles/header.tex — shared LaTeX/Beamer preamble for this project
%
% Use from a Beamer lecture by prepending:
%     \input{header}
% and compiling with TEXINPUTS=../Preambles:$TEXINPUTS (the /compile-latex
% skill does this automatically).
%
% PALETTE CONTRACT. Color names defined here MUST match the names in
% Quarto/theme-template.scss so Beamer and Quarto renderings use the same
% palette. The scripts/check-palette-sync.sh script enforces this.
%
% To customize: edit the HEX values below AND the matching $variable in
% Quarto/theme-template.scss, then run scripts/check-palette-sync.sh.
% =============================================================================

% -----------------------------------------------------------------------------
% Engine detection + encoding
%
% Our /compile-latex skill uses XeLaTeX, where inputenc is unnecessary and
% can produce encoding warnings. Only load inputenc under pdfTeX.
% -----------------------------------------------------------------------------
\usepackage{iftex}
\ifPDFTeX
  \usepackage[utf8]{inputenc}
\fi

\usepackage{amsmath, amssymb, amsthm}
\usepackage{booktabs}
\usepackage{graphicx}

% -----------------------------------------------------------------------------
% PALETTE — mirrors Quarto/theme-template.scss variable names.
% Load xcolor and define colors BEFORE anything that references them
% (notably hyperref, hypersetup, and the Beamer color assignments).
% -----------------------------------------------------------------------------
\usepackage{xcolor}

\definecolor{primary-blue}{HTML}{012169}      % headings, accents
\definecolor{primary-gold}{HTML}{B9975B}      % emphasis, borders
\definecolor{highlight-yellow}{HTML}{F2A900}  % markers, alerts
\definecolor{light-bg}{HTML}{E8EDF5}          % title / section backgrounds
\definecolor{jet}{HTML}{1A1A1A}               % body text

% Semantic colors (used in diagrams and annotations)
\definecolor{positive}{HTML}{15803D}          % good / observed / identified
\definecolor{negative}{HTML}{B91C1C}          % bad / problematic / bias
\definecolor{neutral}{HTML}{525252}           % reference / context

% Highlight accents (match .hi-* classes in the SCSS)
\definecolor{hi-slate}{HTML}{314F4F}
\definecolor{hi-green}{HTML}{2E8B57}
\definecolor{hi-red}{HTML}{C41E3A}

% Semantic aliases so skill templates and snippets can write
% \textcolor{accent}{...} without hard-coding "primary-blue".
\colorlet{accent}{primary-blue}
\colorlet{accent2}{primary-gold}

% -----------------------------------------------------------------------------
% Hyperref — loaded AFTER xcolor + palette so link colors resolve.
% -----------------------------------------------------------------------------
\usepackage{hyperref}
\hypersetup{colorlinks=true, linkcolor=primary-blue, urlcolor=primary-blue,
            citecolor=primary-blue, pdfborder={0 0 0}}

% -----------------------------------------------------------------------------
% Course code — set once per deck (after \input{header}, before \begin{document})
% via \coursecode{CS401}. Shown in the footer on every slide so a deck's course
% is identifiable without opening the title page. Empty by default.
% -----------------------------------------------------------------------------
\makeatletter
\newcommand{\coursecode}[1]{\gdef\@coursecodetext{#1}}
\gdef\@coursecodetext{}
\makeatother

% -----------------------------------------------------------------------------
% Beamer theme assignments (only apply under Beamer).
%
% We detect Beamer via \@ifclassloaded{beamer}, which is the documented,
% non-internal way. This must be wrapped in \makeatletter/\makeatother
% because \@ifclassloaded contains an @-catcode character.
% -----------------------------------------------------------------------------
\makeatletter
\@ifclassloaded{beamer}{%
  \usefonttheme{professionalfonts}
  \setbeamercolor{structure}{fg=primary-blue}
  \setbeamercolor{frametitle}{fg=primary-blue}
  \setbeamercolor{title}{fg=primary-blue}
  \setbeamercolor{subtitle}{fg=primary-gold}
  \setbeamercolor{author}{fg=primary-blue}
  \setbeamercolor{institute}{fg=neutral}
  \setbeamercolor{date}{fg=primary-gold}
  \setbeamercolor{itemize item}{fg=primary-blue}
  \setbeamercolor{itemize subitem}{fg=primary-gold}
  \setbeamercolor{alerted text}{fg=primary-gold}
  \setbeamercolor{block title}{fg=white, bg=primary-blue}
  \setbeamercolor{block body}{bg=light-bg}
  % Semantic block variants: block=definition/concept (blue), exampleblock=
  % worked example (green/positive), alertblock=key takeaway or caution (gold).
  % Keep to <=2 colored boxes per slide across ALL three variants (INV-7).
  \setbeamercolor{block title example}{fg=white, bg=positive}
  \setbeamercolor{block body example}{bg=light-bg}
  \setbeamercolor{block title alerted}{fg=white, bg=primary-gold}
  \setbeamercolor{block body alerted}{bg=light-bg}

  % Minimal footer: course code (left) + page X / N (right), no navigation clutter
  \setbeamertemplate{navigation symbols}{}
  \setbeamertemplate{footline}{%
    \color{neutral}\footnotesize\hspace{0.7em}\@coursecodetext\hfill
    \insertframenumber/\inserttotalframenumber\hspace{0.7em}\vspace{0.3em}}
}{}
\makeatother

% -----------------------------------------------------------------------------
% TikZ setup — libraries the snippet gallery uses
% -----------------------------------------------------------------------------
\usepackage{tikz}
\usetikzlibrary{arrows.meta, positioning, calc, decorations.pathreplacing,
                fit, shapes.geometric, backgrounds}

% Shared TikZ styles — snippets and hand-written diagrams can pick these up
% via `\begin{tikzpicture}[every node/.style={...}]` or by naming specific
% styles (e.g., `\node[dag-node] (X) at (0,0) {X};`).
\tikzset{
  % Node styles for causal DAGs and similar
  dag-node/.style = {
    draw=primary-blue, fill=light-bg, rounded corners=2pt,
    minimum width=1.2cm, minimum height=0.9cm, align=center, font=\small
  },
  decision-node/.style = {
    draw=primary-gold, fill=white, diamond, aspect=1.8,
    minimum width=1.8cm, minimum height=1.2cm, align=center, font=\small,
    inner sep=1pt
  },
  % Edge styles — observed vs counterfactual
  observed-edge/.style = {-{Latex[length=2.5mm]}, thick, draw=jet},
  counterfactual-edge/.style = {-{Latex[length=2.5mm]}, thick, draw=neutral, dashed},
  confound-edge/.style = {-{Latex[length=2.5mm]}, thick, draw=negative, dashed},
  % Data-point styles for plots
  observed-dot/.style = {fill=primary-blue, draw=primary-blue, circle, inner sep=1.2pt},
  counterfactual-dot/.style = {fill=white, draw=neutral, circle, inner sep=1.2pt}
}

% -----------------------------------------------------------------------------
% Convenience macros
% -----------------------------------------------------------------------------
\newcommand{\muted}[1]{\textcolor{neutral}{#1}}
\newcommand{\key}[1]{\textcolor{primary-gold}{\textbf{#1}}}
\newcommand{\good}[1]{\textcolor{positive}{#1}}
\newcommand{\bad}[1]{\textcolor{negative}{#1}}

% Standout frame macro: full-bleed dark background for section transitions.
% Beamer-only (uses \begin{frame}/\setbeamercolor/\usebeamerfont) -- do not
% call from an article-class file (e.g. Notes/<CODE>/*-notes.tex). \muted,
% \key, \good, \bad, and \coursecode above are plain xcolor/gdef and are
% safe to use from either class; verified when Notes/article-class support
% was added.
% Expand: \transitionslide{Your transition title}
\newcommand{\transitionslide}[1]{%
  \begingroup
  \setbeamercolor{background canvas}{bg=primary-blue}
  \begin{frame}[plain]
    \centering
    \vfill
    {\usebeamerfont{title}\color{white}\Large #1\par}
    \vfill
  \end{frame}
  \endgroup
}

% Section divider frame (Beamer-only), an alternative to \transitionslide with a
% darker two-line style: near-black (jet) background, a small white label line
% above a larger gold title, and a thin gold rule along the bottom edge.
% Expand: \sectiondivider{Section 2}{Pointers}
\newcommand{\sectiondivider}[2]{%
  \begingroup
  \setbeamercolor{background canvas}{bg=jet}
  \begin{frame}[plain]
    \vspace*{3.0cm}
    {\color{white}\ttfamily\large #1\par}
    \vspace{0.5cm}
    {\color{primary-gold}\LARGE #2\par}
    \begin{tikzpicture}[remember picture, overlay]
      \draw[primary-gold, line width=2.5pt]
        ($(current page.south west)+(0,0.1)$) -- ($(current page.south east)+(0,0.1)$);
    \end{tikzpicture}
  \end{frame}
  \endgroup
}

% -----------------------------------------------------------------------------
% Channel watermark — "Concepts That Click" (YouTube).
%
% Article-class files (CompetitiveExam Q/A sets, Notes, Assignments) get a
% light diagonal draft watermark on every page plus a centered footer naming
% the channel. Beamer decks get a subtle TikZ background watermark instead
% (the footline already carries the course code). Centralized here so every
% MCQ PDF is branded without per-file edits.
% -----------------------------------------------------------------------------
\makeatletter
\@ifclassloaded{article}{%
  \usepackage{draftwatermark}
  \SetWatermarkText{Concepts That Click}
  \SetWatermarkScale{1}
  \SetWatermarkFontSize{2cm}
  \SetWatermarkLightness{0.86}
  \SetWatermarkAngle{45}
  \usepackage{fancyhdr}
  \pagestyle{fancy}
  \fancyhf{}
  \fancyfoot[C]{\footnotesize\textcolor{neutral}{Concepts That Click~|~YouTube: \href{https://www.youtube.com/@concepts.thatclick}{@concepts.thatclick}}}
  \renewcommand{\headrulewidth}{0pt}
  \renewcommand{\footrulewidth}{0pt}
  \fancypagestyle{plain}{%
    \fancyhf{}%
    \fancyfoot[C]{\footnotesize\textcolor{neutral}{Concepts That Click~|~YouTube: \href{https://www.youtube.com/@concepts.thatclick}{@concepts.thatclick}}}%
    \renewcommand{\headrulewidth}{0pt}%
    \renewcommand{\footrulewidth}{0pt}%
  }
}{}
\@ifclassloaded{beamer}{%
  \setbeamertemplate{background}{%
    \begin{tikzpicture}[remember picture, overlay]
      \node[rotate=45, opacity=0.055, align=center]
        at (current page.center)
        {\Huge\textcolor{neutral}{Concepts That Click}};
    \end{tikzpicture}%
  }
}{}
\makeatother 
\coursecode{JKSSB}
\renewcommand{\thesection}{20.\arabic{section}}
\newtheorem{example}{Example}[section]
\newenvironment{definitionbox}[1]{%
  \par\noindent\textbf{Definition (#1).}\ \itshape
}{\par\vspace{0.3em}}
\title{Windows Explorer --- Detailed Lecture Notes}
\author{JKSSB Computer Awareness}
\date{\today}

\begin{document}
\maketitle

\section{What Windows Explorer is}
\textbf{Windows Explorer} is the file-management program built into the
Windows operating system. Its job is to let the user see and manage
everything stored on the computer: drives, folders, and files. In Windows 10
and Windows 11 the same program is presented to the user as \textbf{File
Explorer}. Internally it is the same program, whose executable file is
\textbf{explorer.exe}. From one Explorer window a user can browse the entire
tree of storage and perform the everyday operations of copying, moving,
renaming, and deleting files.

\begin{definitionbox}{Windows Explorer (File Explorer)}
The file-management program of Windows that displays drives, folders, and
files in a graphical window and lets the user copy, move, rename, and delete
them. It is launched from the executable \texttt{explorer.exe}.
\end{definitionbox}

The main parts of the Explorer window are:
\begin{center}
\begin{tabular}{p{0.30\textwidth}p{0.62\textwidth}}
\toprule
\textbf{Part} & \textbf{What it shows} \\
\midrule
Navigation pane & The tree of drives and folders on the left. \\
Address bar & The full path of the folder that is currently open. \\
File list & The files and subfolders inside the current folder. \\
Ribbon / command bar & The tabs (New, Cut, Copy, Paste, Delete, View) for actions. \\
Search box & Files found inside the current folder matching the typed text. \\
\bottomrule
\end{tabular}
\end{center}

\section{Drives, folders, subfolders, and files}
Storage is organised as a tree. At the top is a \textbf{drive}, identified by
a drive letter. The hard disk is usually \textbf{C:}, and a USB flash drive
might appear as \textbf{E:} or \textbf{F:}. A drive contains
\textbf{folders}. A folder (also called a \textbf{directory}) is a container
that holds files and other folders. A folder placed inside another folder is
a \textbf{subfolder}. A \textbf{file} is a single item of stored data --- a
document, an image, a song, or a program.

\begin{definitionbox}{Drive}
A storage device or partition on which folders and files are kept, named by
a drive letter such as \texttt{C:} or \texttt{E:}.
\end{definitionbox}

\begin{definitionbox}{Folder (directory)}
A named container that holds files and other folders. A folder inside a
folder is called a subfolder.
\end{definitionbox}

Because folders can hold subfolders, and those can hold more subfolders, the
whole arrangement looks like an inverted tree --- the drive is the root, the
folders are branches, and the files are the leaves. The top-most folder of a
drive, directly under the drive letter, is called the \textbf{root folder}.

\section{The file path}
Every file and folder has one exact \textbf{path}: the address that tells
Windows where it lives. The parts of a path are separated by a
\textbf{backslash} (\verb|\|). For example:
\[
\texttt{C:\textbackslash Users\textbackslash Student\textbackslash
Documents\textbackslash report.docx}
\]
Reading left to right, \texttt{C:} is the drive, \texttt{Users},
\texttt{Student}, and \texttt{Documents} are the successive folders, and
\texttt{report.docx} is the file name including its extension. The path is
shown in the Explorer \textbf{address bar}, and the same information is given
in a file's Properties as its \textbf{location}.

\begin{definitionbox}{Path}
The complete address of a file or folder, from the drive letter through the
chain of folders to the item's own name, with parts separated by a backslash.
\end{definitionbox}

\section{File names and extensions}
A file name normally has two parts joined by a dot: the \textbf{name} and the
\textbf{extension}. In \texttt{report.docx}, \texttt{report} is the name and
\texttt{.docx} is the extension. The extension tells Windows which program
should open the file and what kind of data it holds. Common extensions
include \texttt{.txt} for a text file, \texttt{.docx} for a Word document,
\texttt{.xlsx} for an Excel workbook, \texttt{.pptx} for a PowerPoint
presentation, \texttt{.pdf} for a portable document, \texttt{.jpg} for a
photo, and \texttt{.exe} for a program.

\begin{center}
\begin{tabular}{p{0.22\textwidth}p{0.62\textwidth}}
\toprule
\textbf{Extension} & \textbf{Type of file / format} \\
\midrule
\texttt{.txt} & Plain text file \\
\texttt{.docx} & Microsoft Word document \\
\texttt{.xlsx} & Microsoft Excel workbook \\
\texttt{.pptx} & Microsoft PowerPoint presentation \\
\texttt{.pdf} & Portable Document Format \\
\texttt{.jpg} & JPEG image \\
\texttt{.exe} & Executable program \\
\bottomrule
\end{tabular}
\end{center}

\begin{definitionbox}{Extension}
The letters after the last dot in a file name (for example \texttt{.docx})
that indicate the file's type and the program that opens it.
\end{definitionbox}

One common confusion is between an \emph{extension} and a \emph{format}. An
extension is the label at the end of the name, such as \texttt{.png}; the
format is the way the data is actually encoded, such as PNG. The two usually
match, but the words name different things, so they should not be used as
interchangeable (TRM-09, TRAP-34).

\section{Hidden files and attributes}
Every file and folder carries \textbf{attributes}: settings that tell Windows
how to treat it. The common attributes are \textbf{Read-only}, \textbf{Hidden},
\textbf{System}, and \textbf{Archive}. A \textbf{hidden file} or folder has
the \textbf{Hidden} attribute set and is therefore not displayed by default.
Windows hides important system files so that a user does not move or delete
them by mistake. To display hidden items, open the \textbf{View} tab, choose
\textbf{Show}, and then \textbf{Hidden items} (in Windows 10 and 11). In older
versions the setting is reached through \textbf{Folder Options}.

\begin{definitionbox}{File attribute}
A property that controls how Windows treats a file or folder, such as
Read-only, Hidden, System, or Archive.
\end{definitionbox}

\section{Properties}
\textbf{Properties} is the dialog box that displays the details of a selected
file, folder, or drive. It is opened by pressing \textbf{Alt}\,+\,
\textbf{Enter} or by choosing \textbf{right-click $\rightarrow$ Properties}.
For a file, Properties shows the type (format), the size, the date created,
the date modified, the location (path), and the attributes. For a drive, it
shows the total size, the used space, and the free space. Properties is a
read-and-sometimes-edit view: changes to attributes may be made there, but
the file's contents are not shown.

\section{Selecting files}
Before a file can be copied, moved, renamed, or deleted it must be selected.
A single file is selected by clicking it once. A \textbf{contiguous range} of
files is selected by clicking the first file, holding \textbf{Shift}, and
clicking the last file. \textbf{Non-contiguous} files are selected by holding
\textbf{Ctrl} and clicking each file in turn. All files in a folder are
selected with \textbf{Ctrl}\,+\,\textbf{A}.

\begin{center}
\begin{tabular}{p{0.42\textwidth}p{0.50\textwidth}}
\toprule
\textbf{Selection} & \textbf{Method} \\
\midrule
One file & Click once. \\
A block of files in a row & Click the first, hold Shift, click the last. \\
Separate files scattered in the list & Hold Ctrl and click each. \\
Every file in the folder & Ctrl\,+\,A. \\
\bottomrule
\end{tabular}
\end{center}

\section{Copy}
\textbf{Copy} makes a duplicate of the selected item while leaving the
original where it is. The usual method is to press \textbf{Ctrl}\,+\,\textbf{C}
to copy, open the destination folder, and press \textbf{Ctrl}\,+\,\textbf{V}
to paste. The same steps can be done with \textbf{right-click $\rightarrow$
Copy} followed by \textbf{Paste}. Holding \textbf{Ctrl} while dragging an
item also copies it. The key point about copy is that after the operation
there are \emph{two} copies: the original and the new one.

\begin{example}
A student copies \texttt{report.docx} from \texttt{Documents} to a USB drive.
After pasting, the file exists in both places. The copy in
\texttt{Documents} is unchanged, because Copy never removes the original.
\end{example}

\section{Move}
\textbf{Move} shifts the selected item to a new location and removes it from
the old one. The standard way is to \textbf{cut} with
\textbf{Ctrl}\,+\,\textbf{X}, open the destination, and \textbf{paste} with
\textbf{Ctrl}\,+\,\textbf{V}. Cut followed by Paste is therefore the standard
move operation. Alternatively, dragging a file to a folder on the
\textbf{same drive} moves it. Dragging to a folder on a \textbf{different
drive} copies it instead, because Windows assumes you intend to transfer a
copy between separate devices.

\begin{definitionbox}{Move}
Relocating a file or folder from one place to another so that it no longer
exists at the old location. Cut\,+\,Paste performs a move.
\end{definitionbox}

\section{Rename}
\textbf{Rename} changes only the \textbf{name} of a file or folder; its
contents are untouched. The quickest method is to select the item and press
\textbf{F2}, type the new name, and press \textbf{Enter}. The same command is
available by right-clicking the item and choosing \textbf{Rename}. A warning
worth remembering: changing the text after the dot can change the file's
recognised type, so the extension should be left alone unless the type is
genuinely being changed.

\section{Delete and the Recycle Bin}
Pressing the \textbf{Delete} key on a selected item moves it to the
\textbf{Recycle Bin}. The Recycle Bin is a special folder that temporarily
holds deleted items, giving the user a chance to recover them. An item in the
Recycle Bin may be \textbf{restored} to its original location; the
\textbf{Restore} command puts it back. Choosing \textbf{Empty Recycle Bin}
deletes everything held in the Bin permanently.

\begin{definitionbox}{Recycle Bin}
A special Windows folder that temporarily stores files and folders deleted
with the Delete key, so that they can be restored or removed permanently.
\end{definitionbox}

\section{Recycle Bin vs permanent deletion}
An item deleted with \textbf{Delete} is recoverable, because it sits in the
Recycle Bin until the Bin is emptied. Permanent deletion happens when the
user presses \textbf{Shift}\,+\,\textbf{Delete}, when the Recycle Bin is
emptied, or when an item is deleted from a \textbf{USB drive}, a
\textbf{network drive}, or another location that does not use the Recycle
Bin. Once deleted permanently, a file cannot be restored through the normal
Windows interface.

\begin{center}
\begin{tabular}{p{0.30\textwidth}p{0.30\textwidth}p{0.30\textwidth}}
\toprule
\textbf{Point} & \textbf{Recycle Bin (Delete)} & \textbf{Permanent deletion} \\
\midrule
Shortcut & Delete key & Shift\,+\,Delete \\
Recoverable & Yes, can be restored & No, cannot be restored \\
Stored where & Recycle Bin folder & Removed from the disk \\
\bottomrule
\end{tabular}
\end{center}

Permanent deletion does not mean the storage space is instantly overwritten;
the freed space is simply marked as available for reuse. For exam purposes
the distinction is the one that matters: Delete is recoverable, Shift+Delete
is not.

\section{Keyboard shortcuts and drag-and-drop}
The following shortcuts are high-yield for the exam and should be memorised
exactly.

\begin{center}
\begin{tabular}{p{0.40\textwidth}p{0.52\textwidth}}
\toprule
\textbf{Action} & \textbf{Shortcut} \\
\midrule
Copy & Ctrl\,+\,C \\
Cut (move) & Ctrl\,+\,X \\
Paste & Ctrl\,+\,V \\
Undo & Ctrl\,+\,Z \\
Select all & Ctrl\,+\,A \\
Rename & F2 \\
Refresh & F5 \\
Properties & Alt\,+\,Enter \\
New folder & Ctrl\,+\,Shift\,+\,N \\
Delete to Recycle Bin & Delete \\
Permanent delete & Shift\,+\,Delete \\
\bottomrule
\end{tabular}
\end{center}

Drag-and-drop follows a fixed rule. Dragging an item to a folder on the
\textbf{same drive moves} it; dragging to a folder on a \textbf{different
drive copies} it. Holding \textbf{Ctrl} while dragging always \textbf{copies},
and holding \textbf{Shift} while dragging always \textbf{moves}, regardless
of the drives involved.

\section{Exam retrieval and common confusions}
Five distinctions prevent most errors on this topic.
\begin{enumerate}
  \item \textbf{Copy} leaves the original in place; \textbf{move}
    (Cut\,+\,Paste) removes it. Do not confuse the two.
  \item The \textbf{Delete} key sends a file to the \textbf{Recycle Bin}; it
    is not a permanent erase. \textbf{Shift\,+\,Delete} is the permanent
    erase and bypasses the Bin.
  \item Deleting from a \textbf{USB drive} or \textbf{network drive} bypasses
    the Recycle Bin even when the plain Delete key is used.
  \item An \textbf{extension} is the label after the dot (for example
    \texttt{.png}); it is not the same as the \textbf{format} (PNG)
    (TRM-09).
  \item Same-drive drag \textbf{moves}; different-drive drag \textbf{copies}.
    \textbf{Ctrl+drag} forces a copy; \textbf{Shift+drag} forces a move.
\end{enumerate}

\begin{example}
An item asks where a file goes when the Delete key is pressed. Trace the rule:
the Delete key moves the item to the Recycle Bin, from which it can still be
restored. If the stem instead says the file must be removed so that it cannot
be recovered, the required operation is Shift+Delete (or emptying the Recycle
Bin), not the plain Delete key.
\end{example}

\subsection*{Exercises}
\begin{enumerate}[label=\arabic*.]
  \item What is Windows Explorer called in Windows 11, and what is the name
    of its executable program?
  \item Distinguish between a folder and a subfolder, and define a file path.
  \item What is a file extension, and why is an extension not the same as a
    file format?
  \item Explain the difference between Copy and Move, and give the keyboard
    shortcut for each.
  \item A file is deleted with the Delete key. Can it be recovered? Explain,
    and state one way to delete a file permanently.
\end{enumerate}

\subsection*{Solutions}
\begingroup\small
\begin{enumerate}[label=\arabic*.]
  \item In Windows 11 it is called \textbf{File Explorer}. Its executable
    program is \textbf{explorer.exe}.
  \item A folder is a container that holds files and other folders; a folder
    placed inside another folder is a subfolder. A path is the complete
    address of a file or folder, from the drive letter through the chain of
    folders to the item's own name, with parts separated by a backslash.
  \item An extension is the letters after the last dot in a file name (for
    example \texttt{.docx}) that indicate the file's type. It is not the same
    as the format, because the extension is only the label at the end of the
    name, while the format is the way the data is actually encoded.
  \item \textbf{Copy} duplicates the item and leaves the original in place
    (Ctrl\,+\,C, then Ctrl\,+\,V). \textbf{Move} relocates the item and
    removes it from the old location (Ctrl\,+\,X, then Ctrl\,+\,V).
  \item Yes, a file deleted with the Delete key can normally be recovered,
    because it goes to the Recycle Bin and can be restored. It can be deleted
    permanently with \textbf{Shift\,+\,$\boxed{\text{Delete}}$} (or by
    emptying the Recycle Bin).
\end{enumerate}
\endgroup
\end{document}
